\section{The weighted reduction}
\label{sec:construction}

We first construct a digraph with positive rational vertex costs.  Write
$w_v$ for the cost of a vertex $v$ and $w(F)=\sum_{v\in F}w_v$ for the cost
of a vertex set $F$.  For each
fixed integer $A\ge 1$, the construction will have a feedback vertex set of
cost at most $4$ in the completeness case of Theorem~\ref{thm:games}, whereas
every feedback vertex set will have cost greater than $D=10A$ in its soundness
case.  Section~\ref{sec:unweighting} converts this gap into an unweighted one.

\subsection{Choice of parameters}
\label{sec:parameters}

Fix an integer $A\ge 1$ and put $D=10A$. The constants have the following
roles. There will be $M$ bias scales and at most $T$ prefix slots.
The full-prefix and comparison tests have total costs $H$ and $K$.
The junta size is bounded by $J$, independently of the game alphabets;
$\theta$ is the required game accuracy. Finally $L$ and $N$ specify
finite ordered rank sets, and $\psi:[L]\hookrightarrow[N]$ is an
increasing random embedding. Its purpose is to make the order
comparisons insensitive to their precise rank levels, as proved in
Lemma~\ref{lem:rank-spreading}.
Choose the constants in the following order.
\begin{enumerate}[label=\arabic*.]
  \item Set
  \[
    M=128(D+1),\qquad
    p_i=2^{-(i+1)},\qquad q_i=2p_i\quad (i\in[M]),
    \qquad \gamma=\frac{p_M}{1000}.
  \]
  \item Set $H=8(D+1)$.  Choose an integer $T\ge 1$ large enough that
  \begin{equation}
    HM\exp\left(-\frac{p_MT}{8M}\right)\le 1,
    \label{eq:prefix-length-choice}
  \end{equation}
  and an integer
  \begin{equation}
    K\ge \frac{M2^T D}{\gamma}.
    \label{eq:comparison-mass-choice}
  \end{equation}
  \item Put $B_0=(D+1)/\gamma$.  By Corollary~\ref{cor:dyadic-junta}, there
  is an integer $J\ge 1$ that bounds the size of a junta approximant with
  error at most $\gamma$ for every function of total influence at most $B_0$
  at any of the biases $p_i,q_i$, $i\in[M]$.  Choose a positive rational
  $\theta$ small enough that Theorem~\ref{thm:games} applies and
  \begin{equation}
    K\theta\le 1,\qquad J^2\theta\le\gamma.
    \label{eq:game-accuracy-choice}
  \end{equation}
  Apply that theorem with this fixed $\theta$, and let $X,Y$ be its
  alphabets, with $m=\abs{X}=2\abs{Y}$.
  \item Put
  \[
    k_{\max}=2^Tm,\qquad L=10k_{\max}+30,
    \qquad \rho=\min\{\gamma,1/m\}.
  \]
  Choose a positive rational $\eta$ satisfying
  \begin{equation}
    3\eta L^3 2^{m^T}\le\rho.
    \label{eq:rank-accuracy-choice}
  \end{equation}
  Apply Lemma~\ref{lem:rank-spreading} to $L,\eta$, obtaining a finite $N$
  and a rational probability distribution on increasing maps
  $\psi:[L]\hookrightarrow[N]$.  Fix this distribution for the construction.
\end{enumerate}

All these quantities depend only on $A$.  In particular, the junta bound
$J$ is fixed \emph{before} the alphabet size $m$ is determined.  Its existence
does not require an alphabet bound.  The later choices of $\rho,\eta,N$
therefore introduce no dependence back into the choice of $\theta$.

Let the input game have bipartition $U,V$, edge distribution $\nu$, and
endpoint marginals $\nu_U,\nu_V$.  For an edge $e=(u,v)$, write
$\pi_e:X\to Y$ for its two-to-one projection.  All probability distributions
below are finite and rational.  A sampling description specifies a separate
vertex for each outcome of positive probability, with the indicated
probability as a factor in its cost.  The reduction enumerates these outcomes;
it does not sample a random output graph.

\subsection{Rank vertices and their identifications}

\begin{definition}[Rank graph]
\label{def:rank-graph}
For $w\in U\cup V$, let $\Lambda_w=X$ if $w\in U$ and $\Lambda_w=Y$ if
$w\in V$, and put $m_w=\abs{\Lambda_w}$.  The allowed vertex tuples are
\[
  \bigcup_{t=0}^{T}U^t
  \quad\cup\quad
  \bigcup_{t=0}^{T-1}(U^t\times V).
\]
Tuples are ordered sequences and may contain repeated vertices.  For an
allowed tuple $\mathbf d=(w_1,\ldots,w_t)$, let
$\Lambda_{\mathbf d}=\prod_{j=1}^t\Lambda_{w_j}$; the product for the empty
tuple is a singleton.  In every case
\begin{equation}
  \abs{\Lambda_{\mathbf d}}\le m^T.
  \label{eq:tuple-domain-bound}
\end{equation}
For every function $x:\Lambda_{\mathbf d}\to[N]$, introduce a formal rank
vertex $b_{\mathbf d}(x)$.  Within each tuple, introduce an arc
$b_{\mathbf d}(x)\to b_{\mathbf d}(y)$ whenever
\[
  x(\lambda)<y(\lambda)
  \qquad\text{for every }\lambda\in\Lambda_{\mathbf d}.
\]
Whenever $\mathbf s w$ is an allowed extension of $\mathbf s$, identify
\begin{equation}
  b_{\mathbf s}(x)
  \equiv b_{\mathbf s w}(x\circ\operatorname{proj}),
  \label{eq:prefix-identification}
\end{equation}
where $\operatorname{proj}:\Lambda_{\mathbf s w}\to\Lambda_{\mathbf s}$
drops the last coordinate.  Take the equivalence relation generated by all
these identifications and keep the induced arcs between its classes.  Each
resulting rank vertex has cost $D+1$, independently of the number of its
formal representatives.  We use the same notation for a formal vertex and
its equivalence class.  The vertices represented by the constant profiles
with value $j$ are denoted by $o_j$, for $j\in[N]$.
\end{definition}

\begin{lemma}[Evaluation invariant]
\label{lem:evaluation-invariant}
Every global game labeling $\ell$ gives a well-defined real value on the
rank vertices by
\[
  b_{\mathbf d}(x)\longmapsto x(\ell(\mathbf d)),
\]
where $\ell(\mathbf d)$ is the tuple of labels supplied by $\ell$.  Every
rank arc strictly increases this value.  Consequently, the rank graph has
no self-loops, and $o_1,\ldots,o_N$ are distinct vertices with an arc
$o_j\to o_{j'}$ whenever $j<j'$.
\end{lemma}

\begin{proof}
For every generating identification in
Equation~\eqref{eq:prefix-identification},
\[
  (x\circ\operatorname{proj})(\ell(\mathbf s),\ell(w))
  =x(\ell(\mathbf s)).
\]
Thus evaluation is constant on each equivalence class.  This remains true
when vertices repeat in a tuple: their repeated coordinates are evaluated
at the same global label.  A rank arc is strictly increasing at every tuple
of labels, and hence at the particular tuple supplied by $\ell$.  Its two
endpoints cannot become the same equivalence class.  Finally, constant
profiles all identify with their representatives at the empty tuple, and
their evaluated values are the distinct integers $1,\ldots,N$.
\end{proof}

This division into protected rank vertices and deletable tests is related
to the bit-vertex/test-vertex architecture of
\cite[Section~3.1.1]{Svensson2013}; here the rank profiles and prefix
identifications enforce the consistency needed below.
All other vertices will be test vertices, distinct from the rank vertices
and from one another.  Their arcs will be incident only to rank vertices.
For $E\subseteq\Lambda_{\mathbf d}$ and $1\le a<h\le L$, write
\[
  b_{\mathbf d}^{\psi}(E;a,h)=b_{\mathbf d}(x),\qquad
  x(\lambda)=
  \begin{cases}
    \psi(a),&\lambda\in E,\\
    \psi(h),&\lambda\notin E.
  \end{cases}
\]

\subsection{Wildcard vertices and prefix data}

\begin{definition}[Wildcard vertex]
\label{def:wildcard}
Fix an allowed tuple $\mathbf d$, an increasing map $\psi$, an ordered
partition $C_1,\ldots,C_k$ of $\Lambda_{\mathbf d}$ with
$1\le k\le k_{\max}$, and a set $Z\subseteq\Lambda_{\mathbf d}$, called
the star set.  Empty cells in the partition are permitted.  Put
$z_\ell=\psi(10\ell+10)$ for $\ell\in[k]$.  The corresponding wildcard
vertex has an incoming arc from every rank vertex $b_{\mathbf d}(x)$ for
which
\begin{equation}
  x(\lambda)<z_\ell
  \quad\text{for all }\ell\in[k]
  \text{ and }\lambda\in C_\ell\setminus Z,
  \label{eq:wildcard-incoming}
\end{equation}
and an outgoing arc to every $b_{\mathbf d}(x)$ satisfying the strict
reverse inequalities
\begin{equation}
  x(\lambda)>z_\ell
  \quad\text{for all }\ell\in[k]
  \text{ and }\lambda\in C_\ell\setminus Z.
  \label{eq:wildcard-outgoing}
\end{equation}
There are no restrictions on the values of $x$ on $Z$.
\end{definition}

The wildcard inequalities have a direct interpretation under any global
labeling. If its evaluated tuple lies in $C_r\setminus Z$, the threshold
$z_r$ lies strictly between the evaluations of every incoming and every
outgoing rank vertex. Assigning that threshold to the test therefore extends
the increasing evaluation along its incident arcs. At starred tuples, the
arc conditions impose no restriction on these evaluated values; the
completeness argument will delete the corresponding test.

\begin{definition}[Prefix data]
\label{def:prefix-data}
For $t\in\{0,\ldots,T\}$, draw
$\mathbf s=(u_1,\ldots,u_t)$ independently from $\nu_U$ in each slot.
Independently draw $i_1,\ldots,i_t$ uniformly from $[M]$.  Conditional on
these indices, draw $P_j\subseteq X$ using independent Bernoulli
$p_{i_j}$ membership bits, independently for different slots $j$.
Partition $\Lambda_{\mathbf s}$ by the membership patterns
\[
  \bigl(\ind{\lambda_j\in P_j}\bigr)_{j=1}^t,
  \qquad \lambda\in\Lambda_{\mathbf s}.
\]
Order the $2^t$ cells lexicographically with $1$ before $0$, keeping cells
that happen to be empty, and denote them by $C_1,\ldots,C_{2^t}$.  Put
\[
  C_{<r}=\bigcup_{\ell<r}C_\ell,
  \qquad C_{\le r}=\bigcup_{\ell\le r}C_\ell.
\]
At $t=0$ this is the one-cell partition of the singleton product domain.
\end{definition}

\subsection{The three test families}

\begin{definition}[Test families]
\label{def:tests}
Add the following families of test vertices to the rank graph.  In each
family, independently draw $\psi$ from the fixed rank distribution.  Each
outcome is a separate vertex with cost equal to its probability multiplied
by the stated total mass.  Distinct outcomes remain separate even if they
produce the same incident arcs.
\begin{enumerate}[label=\arabic*.,leftmargin=*]
  \item \textbf{Full-prefix tests, of total mass $H$.}
  \label{item:full-prefix-test}
  Draw prefix data at length $T$ and use the wildcard vertex in tuple
  $\mathbf s$ with those $2^T$ cells.  Its star set is
  \begin{equation}
    Z=\left\{\lambda\in\Lambda_{\mathbf s}:
      \text{there exists }i\in[M]\text{ such that }
      \#\{j:i_j=i,\ \lambda_j\in P_j\}
      <\frac{p_iT}{2M}\right\}.
    \label{eq:full-prefix-stars}
  \end{equation}

  \item \textbf{Single-label tests, with one family on each side.}
  \label{item:single-label-test}
  The big-side family has total mass $\abs{X}$, and the small-side family
  has total mass $\abs{Y}$.  Choose $t$ uniformly from
  $\{0,\ldots,T-1\}$ and draw prefix data at length $t$.  Independently
  draw $w$ from the endpoint marginal on the chosen side, a uniformly
  random label $j_0\in\Lambda_w$, and a uniformly random permutation
  $(\sigma(1),\ldots,\sigma(m_w))$ of $\Lambda_w$.  Use a wildcard vertex
  in tuple $\mathbf s w$.  Its cells, in order, are
  \[
    C_\ell\times\{\sigma(a)\},
    \qquad \ell=1,\ldots,2^t,\quad a=1,\ldots,m_w,
  \]
  with the prefix-cell index first.  Thus the same label permutation is
  used in every prefix cell.  Its star set is
  $Z=\Lambda_{\mathbf s}\times\{j_0\}$.

  \item \textbf{Comparison tests, of total mass $K$.}
  \label{item:comparison-test}
  Choose $t$ uniformly from $\{0,\ldots,T-1\}$ and draw prefix data at
  that length.  Independently draw an edge $e=(u,v)$ from $\nu$ and an
  index $i$ uniformly from $[M]$.  Conditional on $t$, choose $r$ uniformly
  from $[2^t]$, independently of all other choices.

  Draw $R_u\subseteq X$ and $R_v\subseteq Y$ by the following coupling.
  Independently in each block indexed by $y\in Y$, draw the two bits on
  $\pi_e^{-1}(y)$ independently with probability $p_i$ of being one.  If
  either is one, set the bit at $y$ to one.  If both are zero, set it to
  one with probability
  \begin{equation}
    \frac{p_i^2}{(1-p_i)^2}.
    \label{eq:coupling-residual}
  \end{equation}
  This is a valid probability because $p_i\le 1/4$.  The marginal
  probability of the bit at $y$ being one is
  \[
    1-(1-p_i)^2+(1-p_i)^2\frac{p_i^2}{(1-p_i)^2}
    =2p_i=q_i.
  \]
  Hence the two complete marginals are the product measures of biases
  $p_i$ and $q_i$, respectively.  Different projection blocks are
  independent, and
  \begin{equation}
    R_u\subseteq\pi_e^{-1}(R_v)
    \label{eq:coupling-inclusion}
  \end{equation}
  holds for every outcome.  These bits are fresh, independent of the
  prefix subsets.

  For $w\in\{u,v\}$ and $R\subseteq\Lambda_w$, define
  \begin{equation}
    E_w^r(R)=
      (C_{<r}\times\Lambda_w)\cup(C_r\times R).
    \label{eq:comparison-sets}
  \end{equation}
  The test vertex has just the incoming and outgoing arcs in the path
  \begin{equation}
    b_{\mathbf s v}^{\psi}(E_v^r(R_v);2,17)
    \ \longrightarrow\ \text{test}\ \longrightarrow\
    b_{\mathbf s u}^{\psi}(E_u^r(R_u);3,18).
    \label{eq:comparison-path}
  \end{equation}
\end{enumerate}
\end{definition}

All wildcard partitions have admissible size: the full-prefix family has
$2^T\le k_{\max}$ cells, and a single-label test has
$2^tm_w\le 2^{T-1}m\le k_{\max}$ cells.  Their tuples, as well as both
tuples in a comparison test, are among those of
Definition~\ref{def:rank-graph}.  The padding in $L$ contains every rank
index used above.

\subsection{Completeness}

\begin{proposition}
\label{prop:completeness}
If the input game has a labeling satisfying a fraction at least $1-\theta$
of its constraints under $\nu$, then the constructed digraph has a
feedback vertex set of cost at most $4$.
\end{proposition}

\begin{proof}
Fix such a labeling $\ell$.  Delete each wildcard test whose evaluated
tuple $\ell(\mathbf d)$ lies in its star set, and each comparison test
whose edge is not satisfied by $\ell$.  Denote this deletion set by
$F_{\mathrm{yes}}$.  No rank vertex is deleted.

We first bound its cost.  In a full-prefix test, conditional on any fixed
vertex tuple $\mathbf s$, the count
\[
  \#\{j:i_j=i,\ \ell(u_j)\in P_j\}
\]
is binomial with mean $\mu_i=p_iT/M$.  Indeed, indices and subsets are
independent between slots, so every slot contributes an independent
Bernoulli variable of parameter $p_i/M$.  This argument also applies
when the same game vertex appears in several slots.  For a binomial random
variable $Z$ of mean $\mu$, the exponential Markov inequality gives
\[
  \Pr[Z<\mu/2]
  \le e^{\mu/2}\E[e^{-Z}]
  \le\exp\bigl(-(1/2-e^{-1})\mu\bigr)
  \le\exp(-\mu/8).
\]
Here the middle inequality follows by writing
$\E[e^{-Z}]=(1-a+ae^{-1})^T\le\exp(-Ta(1-e^{-1}))$
when $Z$ has parameters $T,a$.  Applying this bound to each scale and
taking a union bound over $i\in[M]$, the deletion probability for a
full-prefix test is at most
\[
  \sum_{i=1}^{M}\exp(-\mu_i/8)
  \le M\exp\left(-\frac{p_MT}{8M}\right).
\]
After multiplication by the mass $H$, this is at most one by
Equation~\eqref{eq:prefix-length-choice}.  In a single-label test, membership
of the evaluated tuple in the star set is precisely the event
$j_0=\ell(w)$.  Its probability is $1/m_w$, so each of the two families
contributes cost one.  Comparison tests are deleted with probability at
most $\theta$, and hence contribute at most $K\theta\le1$ by
Equation~\eqref{eq:game-accuracy-choice}.  Thus
\[
  w(F_{\mathrm{yes}})\le 1+1+1+1=4.
\]

To prove that this deletion set is a feedback vertex set, assign to each
rank vertex its evaluation under $\ell$, as in
Lemma~\ref{lem:evaluation-invariant}.  Every rank arc strictly increases
the assigned value.  For a retained wildcard vertex, let $C_r$ be the
cell containing its evaluated tuple and assign the test value $z_r$.
The evaluated tuple is outside the star set.  Thus
Equations~\eqref{eq:wildcard-incoming} and~\eqref{eq:wildcard-outgoing}, at that
tuple, put every incoming rank value strictly below $z_r$ and every
outgoing rank value strictly above it.

For a retained comparison, the evaluated prefix is the same on both
sides and $\pi_e(\ell(u))=\ell(v)$.  By
Equation~\eqref{eq:coupling-inclusion},
\[
  \ell(\mathbf s u)\in E_u^r(R_u)
  \quad\Longrightarrow\quad
  \ell(\mathbf s v)\in E_v^r(R_v).
\]
The possible pairs of rank indices at the start and end of
Equation~\eqref{eq:comparison-path} are therefore $(2,3)$, $(2,18)$, and
$(17,18)$.  Each is increasing, so the corresponding evaluated values are
strictly increasing as well.  Assign this test any real value strictly
between its two endpoint values.

These assignments coexist because there are no arcs between test vertices.
Every arc remaining after deletion of $F_{\mathrm{yes}}$ strictly
increases the assigned real value.  A directed cycle would force a strict
increase back to its starting value, which is impossible.  Thus the
remaining digraph is acyclic.
\end{proof}

\subsection{Deterministic implementation}

\begin{proposition}[Deterministic polynomial construction]
\label{prop:polynomial}
For each fixed $A$, the preceding construction produces a loopless digraph
with positive rational vertex costs in deterministic polynomial time in the
game input size.  Its vertex costs are bounded above by a constant depending
only on $A$.
\end{proposition}

\begin{proof}
Put $n=1+\abs{U}+\abs{V}+\abs{E}$, counting parallel constraints as
distinct edges if necessary.  The number of allowed tuples is at most
$(2T+1)n^T$.  By Equation~\eqref{eq:tuple-domain-bound}, a tuple has at most
$N^{m^T}$ formal rank vertices.  Both factors other than $n^T$ are constants
for fixed $A$.

For each slot, there are at most $M2^m$ choices of index and subset.  The
support of the rank law, the number of permutations of either alphabet,
and all choices of individual labels are also bounded by constants.  A
full-prefix test uses $T$ game-vertex slots, so this family has
$O_A(n^T)$ outcomes.  A single-label test has at most $T-1$ prefix vertices and one
endpoint, again giving $O_A(n^T)$ outcomes.  A comparison test has at most
$T-1$ prefix vertices and one input edge; its remaining choices, including
the coupled subsets and candidate cell, range over constant finite sets.
Thus there are $O_A(n^T)$ comparison outcomes as well.  Here $O_A$ allows
the implicit constant to depend on $A$.

Within a tuple there are at most $N^{2m^T}$ formal rank arcs, and a wildcard
vertex has at most $2N^{m^T}$ incident arcs.  Comparisons have two arcs
each.  Therefore the total numbers of formal vertices and arcs are
$O_A(n^T)$.  Computing the equivalence relation in
Equation~\eqref{eq:prefix-identification} and discarding duplicate arcs
cannot increase these counts.  Lemma~\ref{lem:evaluation-invariant} rules
out rank self-loops, and every other arc joins a rank vertex to a distinct
test vertex.

It remains to account for numerical representation.  Marginal endpoint
probabilities are sums of input rational edge probabilities.  They can be
computed with polynomial bit length: a common denominator can be formed
from the product of the input denominators, whose bit length is at most
their total encoding length.  Every outcome probability is a product of
at most $T+1$ input-dependent edge or endpoint probabilities and a fixed
number of rational factors depending only on $A$.  Its bit length is
therefore polynomial in the input encoding length.  Zero-probability
outcomes are omitted.  The rank law supplied by
Lemma~\ref{lem:rank-spreading} is a fixed finite rational table and may be
hardwired into the reduction for this fixed $A$. Alternatively, a
computable finite Ramsey bound and exhaustive search over finite rational
probability tables recover such a law: each total-variation constraint
is an exact rational inequality, and the strict slack in the proof of
Lemma~\ref{lem:rank-spreading} ensures that the search terminates.
This preprocessing depends only on $A$, not on the input game.

Finally, each outcome probability is at most one.  Every test cost is
therefore at most its family's total mass, while every rank cost is $D+1$.
All costs are bounded by
\[
  C_A=\max\{D+1,H,m,\abs{Y},K\},
\]
which depends only on $A$.  These estimates establish both the claimed
running time and the cost bound.  No efficient dependence on a variable
parameter $A$ is asserted.
\end{proof}
